Language models are increasingly being deployed as autonomous decision-makers in domains where their decisions carry meaningful economic consequences. Recent work studies language models in strategic bargaining environments \citep{araujo2026how}, as tools for synthesizing unstructured borrower information that determines lending outcomes \citep{eisfeldt2024ai}, and as providers of personal financial advice about spending, saving, and portfolio allocation over the life cycle \citep{desilva2026ai}. Implicit in these deployments is the assumption that a model's reasoning capabilities naturally equip it to navigate complex economic tradeoffs over welfare, risk, and resource allocation. However, reasoning capability does not imply economic rationality over such tradeoffs, since the capacity to process complex information is fundamentally distinct from the possession of coherent subjective preferences. Once such systems participate in allocation problems, the relevant questions must begin to address the economic reasoning of language models. We must therefore ask what latent economic preferences govern these models and whether their choices adhere to the axioms of rationality.

In this paper, we develop a structural methodology to quantify the innate economic preferences of language models and formally test their adherence to the rational choice axioms of completeness, monotonicity, and transitivity. Our approach utilizes a formal equivalence between the architecture of language model choice and the Random Utility Model framework. When a model is prompted to select an option from a discrete menu, the raw internal weights it assigns to each alternative before the stochastic realization of text map directly onto the conditional logit choice rule of \cite{mcfadden1972conditional}. This mapping is the foundational identification of the model's latent utility index. Consequently, it places language model choice within the scope of standard structural discrete choice methods. While this measurement framework can accommodate many different economic environments, we deploy a controlled financial portfolio choice setting as our experimental laboratory to test these axioms and estimate structural risk preferences. Each model we test is presented with a menu portfolios that exogenously vary in expected return and risk. Because we can observe the model's latent valuations for each portfolio directly, we can cleanly identify the model's risk tolerance without the confounding noise of stochastic sampling its revealed preferences.

[results paragraph tbd]

% Current evaluation frameworks do not identify the latent preferences that govern language model choice.
% Much of the current literature evaluating language model performance studies them through reasoning benchmarks or through alignment interventions such as instruction tuning, reinforcement learning from human feedback, and direct preference optimization.\footnote{For example, see \cite{hendrycks2021measuring, srivastava2023beyond, ouyang2022training, rafailov2023direct}.} Those exercises reveal a great deal about factual recall, multi-step reasoning, and compliance with preferred rhetoric, but they do not measure the tradeoffs a model is willing to make when faced with competing economic objectives. Success in the prevailing evaluation paradigm can therefore be mistaken for evidence of economic rationality, even though the two objects are conceptually distinct. While logical deduction admits an objective ground truth, economic choice reflects latent preferences over competing tradeoffs. Consequently, attitudes toward risk, time, and inequality must be treated as foundational economic primitives rather than mere cognitive errors to be trained away.

The estimation of structural risk preference parameters from observed choices is already a well-developed program in empirical microeconomics. Field evidence from insurance choices, laboratory evidence from portfolio allocation, and survey-linked experimental evidence have each been used to recover curvature and related preference parameters \citep{cohen2007estimating, choi2007consistency, tanaka2010risk, barseghyan2018estimating}. The language model setting inherits that agenda while altering one feature that matters for identification. Because each model can be estimated separately from a large number of menus, the aggregation problem that motivated random-coefficient and finite-mixture methods for human populations is attenuated \citep{vongaudecker2011heterogeneity, bruhin2010risk}.

First, we define an economically meaningful object for language models in a risky choice environment and show why that object matters, since different risk preferences induce different portfolio rankings and therefore different allocations. Second, we adapt standard structural tools to recover that object from either observed logits or sampled choices. Third, we develop menu-based diagnostics to ask whether the recovered preference is stable enough to support economic interpretation. Fourth, we extend the same measurement framework to training interventions by treating frontier fine tuning as targeted preference induction and by defining the empirical object of interest as the change in the recovered structural preference after training. Taken together, these results imply that agentic AI cannot be deployed in an economically coherent and internally consistent way when the preferences governing its actions are neither measured nor targetable.

Our paper contributes to three literatures. First, it adds to work that studies language models as economic actors by defining a benchmark preference object over risk and return rather than evaluating realized responses alone \citep{horton2023large, chen2023emergence, raman2024steer}. Second, it contributes to the structural literature on risky choice and discrete choice by showing how preference can be recovered from either logits or sampled choices in a common framework \citep{harrison2008risk, choi2007consistency, cohen2007estimating, tanaka2010risk, barseghyan2018estimating, luce1959individual, mcfadden1972conditional}. Third, it contributes to the literature on post-training alignment by treating fine tuning as an intervention on estimated economic preference and by evaluating whether that recovered object is stable enough to support interpretation \citep{christiano2017deep, ouyang2022training, rafailov2023direct, lu2025aligning}. The identified object is deliberately narrow, a benchmark preference parameter recovered from controlled one-token portfolio menus.

The focus is deliberately narrower than the larger question of whether one parameter can summarize all economically relevant model behavior across tasks and domains. What is identified here is a benchmark preference object in one-token portfolio menus, chosen because the environment is clean enough for structural interpretation and because that discipline is needed before broader portability claims can be entertained.

The remainder of the paper proceeds as follows. Section~2 reviews the related literature. Section~3 formalizes language model choice and states the benchmark assumptions. Section~4 develops the portfolio choice laboratory and the mean-variance benchmark. Section~5 describes the experimental design. Section~6 presents the estimation strategy. Section~7 reports the baseline economic preferences. Section~8 studies whether fine tuning can induce targeted preferences. Section~9 discusses economic coherence, internal consistency, and scope. The appendix collects identification results, simulation guidance, and robustness notes.
